\PilotPage{10}

\section*{Harmonic oscillation}

\PilotFigure{p010-harmonic-oscillation}{Harmonic oscillation, peaks, zero crossings, and period.}

The particle vibrates up and down, going through positions $A_0$, $A_1$, $A_2$, $A_3$, $A_4$, $A_5$, and $A_6$. Graphically, this is represented as an oscillatory curve which has peaks and valleys.

The period of oscillation, $\tau$, is the distance in time between two consecutive peaks, $P_1$ and $P_2$, or two consecutive zero crossings with positive slope, $C_1$ and $C_2$.

The frequency, $f$, is how often the particle passes through the same location. Frequency $f$ is the inverse of period $\tau$,
\begin{equation*}
    f=\frac{1}{\tau}\ \mathrm{Hz}.
\end{equation*}

The angular frequency $\omega$ is
\begin{equation*}
    \omega=2\pi f\ \mathrm{rad/s}.
\end{equation*}
